\section*{Capítulo \ref{ch:g12:internal-heat} --- El calor interno de la Tierra}

\begin{solution}{exo:g12:internal-heat:1}
(a) el \omterm{def:g12:internal-heat:heat-flow}{flujo de calor}; (b) la \omterm{def:g12:internal-heat:heat-production}{producción de calor radiogénico}; (c) el
\omterm{def:g12:internal-heat:secular-cooling}{enfriamiento secular}; (d) la \omterm{def:g12:internal-heat:convection}{convección térmica}; (e) la \omterm{def:g12:internal-heat:geothermal}{energía geotérmica}.
\end{solution}

\begin{solution}{exo:g12:internal-heat:2}
(1) \qty{0.0916}{W/m^2}. (2) \qty{4.67e13}{W}. (3) \qty{0.029}{K/m}: una
diferencia de un grado Celsius es una diferencia de un kelvin.
\end{solution}

\begin{solution}{exo:g12:internal-heat:3}
Torio 232, uranio 238, potasio 40 y uranio 235. KamLAND solo ve los
geoneutrinos del torio 232 y del uranio 238.
\end{solution}

\begin{solution}{exo:g12:internal-heat:4}
$490/\sqrt{1} = \qty{490}{mW/m^2}$; $490/2 = \qty{245}{mW/m^2}$;
$490/5 = \qty{98}{mW/m^2}$; $490/7 = \qty{70}{mW/m^2}$, más o menos la media
de los continentes.
\end{solution}

\begin{solution}{exo:g12:internal-heat:5}
Cerca de la superficie, la temperatura del suelo sigue los días y las
estaciones; solo por debajo de unos metros se mantiene en la temperatura
media del año, y solo a más profundidad aún muestra el calor que llega de
abajo. La perforación remueve la roca y el agua del pozo y cambia su
temperatura, así que se espera a que el pozo haya vuelto a la temperatura de
la roca.
\end{solution}

\begin{solution}{exo:g12:internal-heat:6}
(1) $(265 - 10)/9.1 \approx \qty{28.0}{\celsius}$ por kilómetro, es decir,
\qty{0.0280}{K/m}. (2) $q = 2.72 \times 0.0280 \approx \qty{0.0762}{W/m^2}$,
unos \qty{76}{mW/m^2}. (3) Algo por encima de la media de los continentes,
\qty{70.9}{mW/m^2}, alrededor de un \qty{7}{\percent}: un valor continental
típico.
\end{solution}

\begin{solution}{exo:g12:internal-heat:7}
Continentes: $0.0709 \times \num{2.073e14} \approx \qty{1.47e13}{W} =
\qty{14.7}{TW}$. Océanos: $0.1054 \times \num{3.028e14} \approx
\qty{3.19e13}{W} = \qty{31.9}{TW}$. Total, unos \qty{46.6}{TW}, de los que
$31.9/46.6 \approx \qty{68}{\percent}$, dos tercios, salen a través del fondo
oceánico.
\end{solution}

\begin{solution}{exo:g12:internal-heat:8}
(1) $\num{1.88e-6} \times 8000 \approx \qty{0.0150}{W/m^2} = \qty{15.0}{mW/m^2}$.
(2) $15.0/70.9 \approx \qty{21}{\percent}$, una quinta parte del \omterm{def:g12:internal-heat:heat-flow}{flujo de calor}.
\end{solution}

\begin{solution}{exo:g12:internal-heat:9}
$\sqrt{t} = 490/70.9 \approx 6.91$, así que $t \approx 47.8$ millones de
años: el fondo oceánico de más de unos 48 millones de años pierde menos calor
que el continente medio.
\end{solution}

\begin{solution}{exo:g12:internal-heat:10}
Uranio 235: $2^{1.4/0.704} = 2^{1.99} \approx 4$ veces más. Potasio 40:
$2^{2.5/1.248} = 2^{2.00} \approx 4$ veces más. En los dos casos se retroceden
dos periodos de semidesintegración.
\end{solution}

\begin{solution}{exo:g12:internal-heat:11}
(1) $10^2/10^{-6} = \qty{1e8}{s}$, unos 3 años. (2) $(10^5)^2/10^{-6} =
\qty{1e16}{s}$, unos 320 millones de años. (3) $(\num{2.891e6})^2/10^{-6}
\approx \qty{8.4e18}{s}$, unos 265\,000 millones de años. El tiempo crece
como el cuadrado de la distancia: diez veces más grueso, cien veces más
largo.
\end{solution}

\begin{solution}{exo:g12:internal-heat:12}
Multiplica cada uno por $2^{3/T}$: torio 232, $10.4 \times 2^{3/14.0} \approx
\qty{12.1}{TW}$; uranio 238, $9.6 \times 2^{3/4.468} \approx \qty{15.3}{TW}$;
potasio 40, $4.0 \times 2^{3/1.248} \approx \qty{21.2}{TW}$; uranio 235,
$0.42 \times 2^{3/0.704} \approx \qty{8.1}{TW}$. Total, unos \qty{56.7}{TW},
más del doble de los \qty{24.4}{TW} actuales. Los dos isótopos de vida corta
aportan $(21.2 + 8.1)/56.7 \approx \qty{52}{\percent}$, frente a
$4.42/24.4 \approx \qty{18}{\percent}$ hoy.
\end{solution}

\begin{solution}{exo:g12:internal-heat:13}
La potencia es la masa por segundo por el calor específico por el
calentamiento: $\num{3e8} = 1640 \times 4180 \times \Delta T$, así que $\Delta T =
\num{3e8}/(1640 \times 4180) \approx \qty{44}{K}$. El continente medio pierde
\qty{0.0709}{W} por metro cuadrado, así que \qty{300}{MW} son la pérdida
natural de calor de $\num{3e8}/0.0709 \approx \qty{4.2e9}{m^2}$, unos
\qty{4200}{km^2}: la central recoge un calor que la Tierra deja escapar de
forma natural a lo largo de miles de kilómetros cuadrados.
\end{solution}

\begin{solution}{exo:g12:internal-heat:14}
(1) 100 millones de años son $\num{2.6e11}/\num{1e8} = 2600$ veces menos que
260\,000 millones de años; 200 millones de años, 1300 veces menos.
(2) La convección lleva el calor con la roca en movimiento: el tiempo para
atravesar el manto es su espesor dividido por la velocidad, $\qty{2891}{km}$
a unos pocos centímetros al año, unos cien millones de años. La conducción no
tiene ese vehículo: su tiempo crece como el cuadrado de la distancia, y a lo
largo de miles de kilómetros se vuelve enorme.
\end{solution}

\begin{solution}{exo:g12:internal-heat:15}
(1) $10/22.3 \approx \qty{45}{\percent}$; el resto es el enfriamiento del
propio manto, que pierde su propio calor almacenado. (2) El límite está a un radio de $6371 - 2891 =
\qty{3480}{km}$; su superficie es $4\pi (\num{3.48e6})^2 \approx
\qty{1.52e14}{m^2}$, y $\num{1e13}/\num{1.52e14} \approx
\qty{0.066}{W/m^2}$, unos \qty{66}{mW/m^2}, alrededor del \qty{72}{\percent}
de la media en la superficie. (3) El calor atraviesa el límite por conducción
hasta la capa más baja del manto; al calentarse, esta roca se vuelve más
ligera que el manto que tiene encima y puede ascender, llevando el calor
hacia arriba: así se cree que nacen las plumas, en lo profundo.
\end{solution}

\begin{solution}{pb:g12:internal-heat:1}
\textbf{1.} $4\pi (\num{6.371e6})^2 \approx \qty{5.101e14}{m^2}$, y
$2.073 + 3.028 = 5.101$: sí.

\textbf{2.} $3.028/5.101 \approx \qty{59}{\percent}$, unas tres quintas partes.

\textbf{3.} El fondo oceánico es joven: se forma caliente en las dorsales y se
enfría a medida que se aleja, y la \omterm{def:g11:lithosphere-asthenosphere:lithosphere}{litosfera} joven pierde mucho calor.

\textbf{4.} $490/\sqrt{4} = \qty{245}{mW/m^2}$.

\textbf{5.} $2.72 \times 0.026 \approx \qty{0.0707}{W/m^2} =
\qty{70.7}{mW/m^2}$: justo por debajo de la media, prácticamente igual a ella.

\textbf{6.} $0.0709 \times \num{2.073e14} \approx \qty{14.7}{TW}$.

\textbf{7.} $0.1054 \times \num{3.028e14} \approx \qty{31.9}{TW}$.

\textbf{8.} $14.7 + 31.9 \approx \qty{46.6}{TW}$, los \qty{47}{TW} de la
estimación publicada (la media global del capítulo da \qty{46.7}{TW}, el valor
que se usa a continuación).

\textbf{9.} $\num{4.66e13}/\num{5.101e14} \approx \qty{0.0914}{W/m^2}$,
unos \qty{91}{mW/m^2}.

\textbf{10.} $0.0243 \times \num{2.073e14} \approx \qty{5.0}{TW}$, un tercio de
los \qty{14.7}{TW} que se pierden a través de los continentes: la
radiactividad de la \omterm{def:g12:continental-crust:continental-crust}{corteza continental} es una gran parte de su \omterm{def:g12:internal-heat:heat-flow}{flujo de
calor}.

\textbf{11.} Por cada kilogramo de uranio, la Tierra contiene \qty{3.9}{kg}
de torio, que liberan $3.9 \times 26.3 \approx \qty{102.6}{\micro W}$, frente
a los \qty{94.4}{\micro W} del uranio 238. El uranio 238 aporta, pues,
$94.4/(94.4 + 102.6) \approx 0.479$ de los \qty{20.0}{TW}: unos
\qty{9.6}{TW}. Uranio 235: $9.59 \times 4.10/94.4 \approx \qty{0.42}{TW}$.

\textbf{12.} $20.0 + 0.42 + 4 \approx \qty{24.4}{TW}$.

\textbf{13.} El granito y las demás rocas de la \omterm{def:g12:continental-crust:continental-crust}{corteza continental}, unas
veinte veces más ricas que las rocas ultramáficas del manto.

\textbf{14.} Mayor: cada isótopo era entonces más abundante, el doble por cada
periodo de semidesintegración hacia atrás; hace 3000 millones de años, el
total era más del doble del actual.

\textbf{15.} $46.7 - 24.4 = \qty{22.3}{TW}$.

\textbf{16.} Del calor almacenado desde que se formó la Tierra: el
\omterm{def:g12:internal-heat:secular-cooling}{enfriamiento secular} del manto y del núcleo.

\textbf{17.} $22.3/46.7 \approx \qty{48}{\percent}$.

\textbf{18.} Del $11/46.7 \approx \qty{24}{\percent}$ al $38/46.7 \approx
\qty{81}{\percent}$: los geoneutrinos estrechan mucho este intervalo.

\textbf{19.} La \omterm{def:g12:internal-heat:convection}{convección del manto}, que transporta la mayor parte del calor;
solo las últimas decenas de kilómetros se atraviesan por conducción.

\textbf{20.} $24.4/46.7 \approx 0.52$: la radiactividad aporta hoy
\textbf{alrededor de la mitad} de la pérdida de calor de la Tierra, y el lento
enfriamiento del planeta, la otra mitad.
\end{solution}
